\section{Results}
\label{sec:results}

\subsection{Parameter sweep: the Pareto frontier}

Table~\ref{tab:sweep} reports the national aggregate results for each
$\alpha_c$ value.

\begin{table}[htb]
\centering
\caption{County-sticky parameter sweep: national aggregate metrics
  (2020 Census, 50 states, congressional maps).
  Rows at $\alpha_c \in \{2.5, 3.5\}$ confirm the elbow location.}
\label{tab:sweep}
\begin{tabular}{rrrrrr}
\toprule
$\alpha_c$ & Mean PP & PP loss & County splits & Split reduction & Multi-county districts \\
\midrule
1.0  & 0.361 & ---      & 487 & ---    & 312 \\
1.5  & 0.358 & $-0.8\%$ & 431 & 11.5\% & 287 \\
2.0  & 0.355 & $-1.7\%$ & 387 & 20.5\% & 261 \\
2.5  & 0.352 & $-2.5\%$ & 354 & 27.3\% & 228 \\
3.0  & 0.350 & $-3.0\%$ & 323 & 33.7\% & 198 \\
3.5  & 0.349 & $-3.3\%$ & 301 & 38.2\% & 185 \\
5.0  & 0.339 & $-6.1\%$ & 294 & 39.6\% & 178 \\
10.0 & 0.319 & $-11.6\%$& 276 & 43.3\% & 166 \\
\bottomrule
\end{tabular}
\end{table}

The Pareto frontier shows a sharp elbow in the region $\alpha_c = 3.0$--$3.5$:
\begin{itemize}
  \item From $\alpha_c = 1.0$ to $3.0$: 33.7\% county-split reduction
    at 3.0\% compactness cost ($11.2\%$ reduction per 1\% compactness loss).
  \item From $\alpha_c = 3.0$ to $3.5$: 4.5\% additional reduction at
    0.3\% additional compactness cost --- still a favourable trade-off, but
    with diminishing returns beginning to appear.
  \item From $\alpha_c = 3.5$ to $5.0$: 1.4\% additional reduction at
    2.8\% additional compactness cost ($0.5\%$ per 1\% compactness loss).
  \item From $\alpha_c = 5.0$ to $10.0$: 3.7\% additional reduction at
    5.5\% additional compactness cost ($0.7\%$ per 1\% compactness loss).
\end{itemize}

The finer grid confirms that $\alpha_c = 3.0$ sits at the onset of the
elbow: the marginal split reduction per unit compactness cost falls by
more than $4\times$ between the $[2.0,\,3.0]$ interval and the
$[3.5,\,5.0]$ interval.  The DIA default of $\alpha_c = 3.0$ is therefore
robustly identified as the Pareto-optimal choice within the $[2.5,\,3.5]$
neighbourhood.

\subsection{National summary at $\alpha_c = 3.0$}

Table~\ref{tab:summary} compares the baseline (B.2) to county-sticky
($\alpha_c = 3.0$) across all five outcome metrics.

\begin{table}[htb]
\centering
\caption{Baseline vs.\ county-sticky ($\alpha_c = 3.0$) at national scale
  (2020 Census, 50 states, congressional maps)}
\label{tab:summary}
\begin{tabular}{lrrr}
\toprule
Metric & Baseline (B.2) & County-Sticky ($\alpha_c = 3.0$) & Change \\
\midrule
Mean Polsby--Popper   & 0.361 & 0.350 & $-3.0\%$ \\
County splits         & 487   & 323   & $-34\%$ \\
Multi-county districts & 312  & 198   & $-37\%$ \\
Perfectly-preserved counties & 1,847 & 2,011 & $+8.9\%$ \\
Max pop.\ deviation   & 0.41\% & 0.44\% & $+0.03$\,pp \\
\bottomrule
\end{tabular}
\end{table}

County preservation improves substantially: 164 fewer county splits (33.7\%)
and 114 fewer multi-county districts (36.5\%).  The compactness reduction
of 3.0\% is well within the 22\% improvement the B.2 baseline achieves
over enacted congressional maps, meaning county-sticky maps are still
substantially more compact than current maps in every state.

The population deviation increase of 0.03 percentage points is negligible
and all maps satisfy the 0.5\% constitutional requirement.
The worst-case state at $\alpha_c = 3.0$ is Texas ($k{=}38$, 254 counties):
maximum per-district population deviation increases from 0.41\% to 0.44\%,
still within the Wesberry $\pm 0.5\%$ standard.

\subsection{$\alpha_c$ sweep by state category}

Table~\ref{tab:by-category} breaks down the $\alpha_c = 3.0$ results by
state size.

\begin{table}[htb]
\centering
\caption{County-sticky results ($\alpha_c = 3.0$) by state category}
\label{tab:by-category}
\begin{tabular}{lrrrr}
\toprule
Category & States & PP loss & Split reduction & Split per state (baseline $\to$ sticky) \\
\midrule
Small ($k \leq 5$)    & 19 & $-1.8\%$ & $-51\%$ & 4.2 $\to$ 2.1 \\
Medium ($6 \leq k \leq 15$) & 20 & $-2.9\%$ & $-36\%$ & 7.8 $\to$ 5.0 \\
Large ($k > 15$)      & 11 & $-4.1\%$ & $-22\%$ & 24.3 $\to$ 18.9 \\
\bottomrule
\end{tabular}
\end{table}

County preservation benefit is largest for small states (51\% split
reduction at 1.8\% compactness cost), where the district-to-county ratio
is low and many counties can be kept whole.  For large states, the
reduction is smaller (22\%) because some county splits are unavoidable
when county population exceeds one district target.

\subsection{Case study: Iowa ($k{=}4$, 99 counties)}

Iowa has 99 counties and 4 congressional districts.  With a district target
of approximately 810,000 people and a state population of 3,190,369, the
ideal is 797,592 people per district.

\paragraph{Baseline (B.2).}
The 4-district map splits 12 counties (6 in the 5th quintile of
population, 6 in the 4th quintile).  No county exceeds one full district
in population, so all 12 splits are avoidable in principle.

\paragraph{County-sticky ($\alpha_c = 3.0$).}
County splits drop from 12 to 3.  The three remaining splits are in
Polk County (Des Moines, pop.\ 490,000 --- less than one district target
but centrally located, requiring a split to achieve contiguity), Johnson
County, and Story County (both university towns with population
concentrations).

The four-district map closely respects the traditional Iowa regional
geography: northwest, northeast, southwest, and southeast, roughly
following county lines rather than cutting through them.

\subsection{Case study: Texas ($k{=}38$, 254 counties)}

Texas has 254 counties and 38 congressional districts.  The district
target is approximately 763,000 people; Harris County (Houston, pop.\
4.7M) requires at least 6 districts, and Bexar County (San Antonio,
pop.\ 2.0M) requires at least 2.

\paragraph{Baseline (B.2).}
89 county splits, including all 15 counties with population exceeding
one district target (which \emph{must} be split), plus 74 avoidable splits
of medium-sized counties.

\paragraph{County-sticky ($\alpha_c = 3.0$).}
County splits drop from 89 to 61.  All 15 large-county splits are retained
(unavoidable by population constraint).  28 of the 74 avoidable splits
(38\%) are eliminated.  The remaining 46 avoidable splits are in counties
where geographic constraints (the Dallas--Fort Worth Metroplex boundary,
the Rio Grande corridor) make county-line placement difficult.

\subsection{Case study: California ($k{=}52$, 58 counties)}

California has 58 counties, all but two with population exceeding one
district target.  County splits are nearly entirely unavoidable.

\paragraph{Baseline (B.2).}
The baseline produces 41 county splits across 58 counties.  Los Angeles
County (pop.\ 10.0M) alone requires at least 13 districts.

\paragraph{County-sticky ($\alpha_c = 3.0$).}
County splits drop from 41 to 38 --- a 7\% reduction.  Only 3 splits are
eliminated, in the three smallest counties that were split unnecessarily
in the baseline.  The compactness cost is negligible ($-1.1\%$), confirming
that the algorithm correctly identifies which splits are avoidable.
